\section{Governing equations}\label{sec:governing}

Let \(x_i\) be Cartesian coordinates and let \(u_i(x_1,x_2,x_3,t)\),
\(\phi(x_1,x_2,x_3,t)\), and \(n(x_1,x_2,x_3,t)\) be the displacement, electric
potential, and mobile-carrier perturbation. The base concentration is \(n_0\) and
the perturbation is \(\Delta n=n-n_0\). The small strain is
\[
\varepsilon_{ij}=\tfrac12(u_{i,j}+u_{j,i}),
\qquad
\varepsilon_{ij,l}=\partial_l\varepsilon_{ij}.
\]

For a \textit{6mm} crystal with the polar \(c\)-axis along \(x_3\), the ordinary
elastic, piezoelectric, and dielectric sectors have the standard Voigt forms
\[
C=\begin{pmatrix}
C_{11}&C_{12}&C_{13}&0&0&0\\
C_{12}&C_{11}&C_{13}&0&0&0\\
C_{13}&C_{13}&C_{33}&0&0&0\\
0&0&0&C_{44}&0&0\\
0&0&0&0&C_{44}&0\\
0&0&0&0&0&(C_{11}-C_{12})/2
\end{pmatrix},\quad
e=\begin{pmatrix}
0&0&0&0&e_{15}&0\\
0&0&0&e_{15}&0&0\\
e_{31}&e_{31}&e_{33}&0&0&0
\end{pmatrix},\quad
\kappa=\operatorname{diag}(\kappa_{11},\kappa_{11},\kappa_{33}).
\]
The strain-gradient tensor is
\[
A_{ijklmn}=\ell^2C_{ijkl}\delta_{mn},
\]
with \(\ell>0\). The electrochemical potential auxiliary is eliminated rather than
treated as an independent physical field.

The retained internal energy density is
\[
\begin{aligned}
w=&\tfrac12C_{ijkl}\varepsilon_{ij}\varepsilon_{kl}
+\tfrac12A_{ijklmn}\varepsilon_{ij,l}\varepsilon_{kl,m}
+e_{kij}\varepsilon_{ij}\phi_{,k}
-\tfrac12\kappa_{kl}\phi_{,k}\phi_{,l}\\
&+\Delta n\,H_{jkl}\varepsilon_{jk,l}.
\end{aligned}
\]
The last term is the only additional interaction. Its variation gives the reciprocal
stress and carrier terms
\[
\Delta\tau_{jkl}=H_{jkl}\Delta n,
\qquad
\Delta\widetilde{\mu}=H_{jkl}\varepsilon_{jk,l},
\]
so that
\[
\sigma_{ij}=C_{ijkl}\varepsilon_{kl}
-\partial_l\tau_{ijkl}
+e_{kij}\phi_{,k},
\qquad
\tau_{ijkl}=A_{ijklmn}\varepsilon_{mn,l}+\Delta\tau_{ijkl}.
\]
The electric displacement is
\[
D_i=e_{ijk}\varepsilon_{jk}-\kappa_{ij}\phi_{,j}.
\]

The quasistatic field equations are
\[
\sigma_{ij,j}-\rho u_i^{,tt}=0,
\qquad
D_{i,i}+q\Delta n=0,
\]
and the carrier field obeys the drift--diffusion balance
\[
\frac{\partial\Delta n}{\partial t}
+\frac1q\nabla\!\cdot\!\mathbf J
=0,
\qquad
\mathbf J=M n\left(-\nabla\widetilde{\mu}\right),
\qquad
\widetilde{\mu}=\frac{q\phi}{k_BT}+\frac{\Delta n}{n_0}+\Delta\widetilde{\mu}.
\]
For harmonic fields \(u_i=\Re\{\hat u_i\exp[i(k_xx+k_yy-\omega t)]\}\), the
tangential wavenumber \(k_y\) is conserved at the edge. The present three-displacement
reduction is a 2.5D Fourier-reduced geometry with a propagation-plane tangential
component, not strict sagittal plane strain.

Dimensional consistency requires \(C_{ijkl}\) in N/m\(^2\), \(A_{ijklmn}\) in
N/m\(^2\) (equivalently J/m\(^3\) for the strain-gradient energy), \(e_{ijk}\) in
C/m\(^2\), \(\kappa_{ij}\) in F/m, \(M\) in m\(^2\)/(V\,s), \(D\) in m\(^2\)/s,
\(n_0\) in m\(^{-3}\), and \(H_{jkl}\) in J/m. These are the units used by the
numerical implementation and by \cref{tab:parameters}.

For the synthetic tests below, \(A=\ell^2C\otimes I\), \(\ell=50\,\)nm, and only
\(H_{333}\) is retained. The \(H\) magnitude is normalized by a strict full-3D
Schur-stability bound rather than by a measured ZnO tensor. The flexoelectric
coefficient \(f\) is set to zero in the computed cases; this is an explicitly
labelled classical-limit control, not a measured vanishing flexoelectric response.
The retained physical constants and their provenance are listed in
\cref{tab:parameters}.
